\section{A valued lift and a submean inequality}\label{sec:valued}

In the synchronization argument, one parameter can approach several
different values. We need to compare the rate at which the tensor
price falls near each of them. The scalar model is a polynomial with
several roots. The valuation records how closely the parameter
approaches a root, and the multiplicity gives the rate of decrease
of the scalar price. For tensors, Lemma~\ref{lem:star} bounds the sum
of these rates by the growth of the price on an enclosing disk.
The scalar comparison is familiar from the potential theory of the
nonarchimedean projective line~\cite{BakerRumely}. Our tensor
inequality follows by interpolation once we construct a price over
a field with a valuation.

\subsection{The Gauss lift}

Let $H$ be an ordinary price of scale $D>0$ over a field $E$, and put
$\phi(A)=\exp(H(A)/D)$, with $\phi(0)=0$. The ring $E[t^{\mathbb Q}]$
consists of finite sums of rational powers of $t$. On its fraction
field $K$ take the additive valuation $\nu$ given by the least
exponent on nonzero elements of the ring. Thus $\nu(t^r)=r$, and a
nonzero constant in $E$ has valuation zero. Write $[t^j]T$ for the
coefficient tensor of $t^j$. A monomial tensor $t^j A$ should have
price $H(A)-j$. For a general tensor we form powers using the same
$t$ in every factor and take the largest coefficient price with
this valuation adjustment.

\begin{theorem}[Gauss lift]\label{thm:gauss}
There is a valued price $p$ of scale $D$ on $K$ extending $H$. On tensors
with entries in $E[t^{\mathbb Q}]$ it is given by
\begin{equation}\label{eq:gauss-price}
 p(T)=\lim_{n\to\infty}\frac1n
       \max_{j\in\mathbb Q}\bigl(H([t^j]T^{\otimes n})-j\bigr).
\end{equation}
\end{theorem}

\begin{proof}
Choose a common denominator $m$ for the exponents and write
$x=t^{1/m}$. Temporarily assign $x$ an arbitrary weight $b>0$. For
a nonzero Laurent tensor $A(x)$ consider the weighted coefficient sum
and its regularization
\[
 N_b(A)=\sum_i e^{-bi/D}\phi([x^i]A),
 \qquad
 p_b(A)=\lim_{n\to\infty}\frac Dn\log N_b(A^{\otimes n}).
\]
The sum bound for $\phi$ makes $N_b$ subadditive. Applying the same
bound to coefficient convolution and using multiplicativity of $\phi$
shows that $N_b$ is submultiplicative. Hence the limit exists by
subadditivity of the logarithms. The lowest coefficient of
$A^{\otimes n}$ is the $n$-th power of the lowest coefficient of $A$,
which gives a finite lower bound for the normalized logarithms.

There are $O_A(n+1)$ exponents in $A^{\otimes n}$. Replacing the sum
by its largest term therefore costs $O_A(\log(n+1)/n)$ in normalized
logarithmic price. It follows that $p_b(A)$ is a limit of maxima of
affine functions of $b$, whose slopes lie in a fixed bounded
interval. In particular, $p_b(A)$ is convex and Lipschitz. Setting
$b=1/m$ gives Equation~\eqref{eq:gauss-price}, independently of the
common denominator.

\smallskip\noindent\emph{Separating the maximal coefficients.}
Submultiplicativity gives $p_b(T\otimes U)\le p_b(T)+p_b(U)$.
For the reverse inequality, a maximal coefficient from each power
may contribute to a convolution coefficient in which other terms
cancel it. We isolate that pair by varying $x$ in the first factor
alone. Multiplying by $x^a$ shifts $p_b$ by $-ba$, so it suffices to
consider nonzero polynomials $T,U$.

First fix $b$ at which $b\mapsto p_b(T)$ is differentiable. For each
$n$, choose indices $i=i_n$ and $j=j_n$ maximizing the weighted
coefficient sizes of $T^{\otimes n}$ and $U^{\otimes n}$. An index
a fixed fraction of $n$ away from $i$ incurs a uniform deficit.
Indeed, for each $\xi>0$ there is a $\delta>0$ such that for all
sufficiently large $n$,
\begin{equation}\label{eq:gauss-deficit}
 \frac{H([x^h]T^{\otimes n})-bh}{n}
   \le p_b(T)-\delta \qquad (|h-i|\ge \xi n).
\end{equation}
Along any asymptotically maximizing sequence with $h/n\to v$,
comparison at $b'$ gives $p_{b'}(T)\ge p_b(T)-(b'-b)v$.
Differentiability forces $v=-\frac{d}{db}p_b(T)$. Compactness of the
interval of possible $h/n$ gives both convergence of $i/n$ to this
number and the deficit in Equation~\eqref{eq:gauss-deficit}.

At $k=i+j$, introduce a variable $w$ and write
\[
\begin{aligned}
 C_s
 &= [x^k w^s]\bigl(T((1+w)x)^{\otimes n}\otimes U(x)^{\otimes n}\bigr)\\
 &= \sum_h \binom hs\, [x^h]T^{\otimes n}\otimes[x^{k-h}]U^{\otimes n}.
\end{aligned}
\]
Changing $s$ changes the integer multipliers of the same coefficient
products. Consecutive values of $s$ suffice to isolate any one
index in a short interval in every characteristic.

Fix $0<\eta<1/4$. Let $I_n$ consist of the nonnegative indices within
$\eta n/3$ of $i$, and put $m_n=|I_n|$. For large $n$,
$m_n-1\le\eta n$. The matrix
$(\binom hs)_{0\le s<m_n,\ h\in I_n}$ has determinant one over
the integers. Indeed, its columns are the coefficients of consecutive
powers of $1+w$ modulo $w^{m_n}$. Multiplication by the first power is
unitriangular, and the powers with exponents $0,\ldots,m_n-1$ form
a unitriangular basis. Inverting this matrix over $E$ gives a linear
combination $C$ of $C_0,\ldots,C_{m_n-1}$ with multiplier one at $i$
and zero on $I_n\setminus\{i\}$. The isolated product has weighted
size $\exp(n(p_b(T)+p_b(U)+o(1))/D)$. By
Equation~\eqref{eq:gauss-deficit} with $\xi=\eta/3$, each of the
$O(n)$ remaining products has weighted size smaller by an exponential factor.
Their scalar multipliers in $E$ have no cost. Applying subadditivity
of $\phi$ to the isolated product, written as $C$ minus these
remaining terms, gives
\begin{equation}\label{eq:gauss-lower}
 e^{-bk/D}\phi(C)\ \ge\
 \exp\!\left(\frac nD\bigl(p_b(T)+p_b(U)-o(1)\bigr)\right).
\end{equation}

To bound the same combination from above, put
$B_\ell(x)=[w^\ell](T((1+w)x)\otimes U(x))$ and $d=\deg T$.
Then $B_0=T\otimes U$, and the term contributing to $C_s$ with
shift degrees $\ell_1,\ldots,\ell_n$ has at most $s$ nonzero
$\ell_a$. If there are $l$ of them, its weighted coefficient size
is at most
\[
 N_b(B_0^{\otimes(n-l)})\,\max(1,N_b(B_1),\ldots,N_b(B_d))^l,
\]
where the maximum is $1$ if $d=0$. Here we have permuted copies
simultaneously at every site to group the unchanged factors.
The number of choices for $l\le\eta n$ is at most
$\sum_{l\le\eta n}\binom nl(1+d)^l$. Its logarithm is
$O_d(n\eta\log(1/\eta)+\log(n+1))$. Moreover
$D\log N_b(B_0^{\otimes(n-l)})=(n-l)p_b(B_0)+o(n)$ uniformly in
this range. The $O(n)$ coefficients used to form $C$ have logarithmic
cost $O(\log(n+1))$. Taking logarithms in Equation~\eqref{eq:gauss-lower},
dividing by $n/D$ and letting $n\to\infty$ gives
\[
 p_b(T)+p_b(U)\le p_b(T\otimes U)+O_{b,T,U}\bigl(\eta\log(1/\eta)\bigr).
\]
Let $\eta\downarrow0$. Differentiability points are dense, so continuity
establishes the identity for every $b>0$.

\smallskip\noindent\emph{Scalars and integral maps.}
For a nonzero scalar Laurent polynomial the lowest coefficient of
each power has the largest weighted size, since all nonzero scalar
coefficients have ordinary price zero. Thus $p(f)=-\nu(f)$.
Multiplicativity now gives a well-defined extension to fractions by
$p(A/f)=p(A)+\nu(f)$.

Consider a local map with polynomial entries of nonnegative degree.
After taking its $n$-th power, the matrix at each site has degree
$O(n)$. Each output coefficient of degree $j$ is a sum over the
degrees chosen at the finitely many sites. There are polynomially
many terms in $n$, and each is an ordinary restriction of an input
coefficient of degree at most $j$. The input coefficient controls
its term by monotonicity of $H$, and the degree weight only decreases.
The resulting logarithmic loss is $O(\log(n+1))$, which vanishes
after normalization. For a matrix of integral rational entries,
choose a common denominator and remove its lowest power from both
numerator and denominator. The numerator matrix then has only
nonnegative powers and the denominator has valuation zero. Doing
this at each site proves monotonicity for every integral restriction.

For the ordinary sum bound, the binomial expansion, with local
permutations of copies, gives
\[
 N_b((T+U)^{\otimes n})\le\sum_{j=0}^n\binom nj
        N_b(T^{\otimes j})N_b(U^{\otimes(n-j)}).
\]
For a fixed nonzero $T$ and all $0\le j\le n$ we have
$D\log N_b(T^{\otimes j})=jp_b(T)+o(n)$ uniformly. The binomial
formula yields $\exp(p(T+U)/D)\le\exp(p(T)/D)+\exp(p(U)/D)$.
Common denominators give this inequality on $K$. The lift agrees with
$H$ on constants, including every diagonal tensor, so the diagonal test of Section~\ref{sec:prices} gives
the direct-sum identity. Set $p(0)=-\infty$.
\end{proof}

\subsection{Prices on disks}

We now need a generic value of the price on each disk of parameters.
Take an algebraically closed field $F$ and an ordinary price on a fixed
finite set of sites over $F$. Enlarge the coefficient field to
$E=F(\mathcal U)$ for an uncountable set $\mathcal U$ of independent
indeterminates. For a rational tensor depending on finitely many of
these variables, define its ordinary price over $E$ as the supremum
over its $F$-valued specializations on the domain of definition.

The fixed-format bound in Equation~\eqref{eq:prices-range} makes this
supremum finite for every nonzero rational tensor. Its fixed-format
price sublevel sets over $F$ are Zariski closed
(Section~\ref{sec:prices-background}). Hence the specializations
within any positive tolerance of the supremum contain a nonempty
open set. Finitely many such sets meet on the common domain of
definition. Applying the price axioms at these simultaneous
specializations and letting the tolerance tend to zero gives the
axioms over $E$. The extension is consistent on subfields generated
by subsets of $\mathcal U$ and invariant under permutations of
$\mathcal U$.

Apply Theorem~\ref{thm:gauss} over $E$. Write $\mathcal O$ for the
valuation ring of $K=\operatorname{Frac}(E[t^{\mathbb Q}])$.
A \emph{fresh} $u\in\mathcal U$ occurs in none of the data being
compared. For $r\in\mathbb Q$ the disk with center $a\in K$ and
valuation bound $r$ is $a+t^r\mathcal O$. Its generic price is the
price at $a+t^r u$. Increasing $r$ shrinks the disk.

\begin{lemma}[Specialization and the star inequality]\label{lem:star}
For a polynomial tensor $T(c)$ over $K$, let $P(c)=p(T(c))$ on the
base sites, or let $P(c)=\beta_p(T(c))$ when $T$ has additional
sites. With a fresh $u$ and $P_*=P(u)$, one has
\[
 P(c)\le P_* \quad (c\in\mathcal O), \qquad P_*=\sup_{a\in F}P(a).
\]
The generic price stays the same after an affine change of variable
$c\mapsto a+bc$ with $a\in\mathcal O$ and $b\in\mathcal O^\times$.
Here $u$ is fresh also for $a$ and $b$.

Suppose $T$ is nonzero of degree $d$. Choose finitely many centers
$c_j\in\mathcal O$ with distinct residues and positive rational
numbers $r_j,r_\infty$. Let
\[
 P_j=P(c_j+t^{r_j}u),\qquad P_\infty=P(t^{-r_\infty}u).
\]
Then
\begin{equation}\label{eq:valued-star}
 \sum_j\frac{P_*-P_j}{r_j}
 \ \le\ \frac{P_\infty-P_*}{r_\infty}\ \le\ d .
\end{equation}
Equivalently, the first inequality bounds the center price $P_*$
by the weighted average of the prices at the ends of the star.
The weights are proportional to $1/r_j$ and $1/r_\infty$. We call
the values on the star \emph{harmonic} if the center equals this
average. The weighted-average inequality also holds for the zero
family.
\end{lemma}

For instance, a scalar polynomial $T(c)=\prod_j(c-c_j)^{m_j}$
with distinct $c_j\in F$ has generic price zero on the unit disk.
On the inner disk at $c_j$ its price is $-m_jr_j$, giving a drop
of $m_j$ per unit of valuation length. On the outer disk its price
is $r_\infty\sum_j m_j$. Both inequalities in
Equation~\eqref{eq:valued-star} are equalities.

\begin{proof}
The zero family is immediate. First let $P(c)=p(T(c))$ with
$T\ne0$, and write $Q(c)=T(c)^{\otimes n}$. Its degree is at most
$nd$. Interpolation on $nd+1$ fresh indeterminates uses only integral
Lagrange weights at any $c\in\mathcal O$, since the nodes have unit
differences. All the samples have price $nP_*$ by permutation symmetry,
so the ordinary sum bound gives
\[
 p(Q(c))\le nP_*+D\log(nd+1).
\]
Dividing by $n$ and letting $n\to\infty$ proves specialization.
Conversely, interpolate at $nd+1$ distinct nodes of $F$ and evaluate
at $u$. The same argument gives $P_*\le\sup_{a\in F}P(a)$. Applying
specialization to an integral affine change of variable and its
integral inverse proves the asserted invariance. Specialization on
nested disks also gives $P_j\le P_*\le P_\infty$.

The inverse Vandermonde matrix at these nodes has integral entries,
so interpolation controls coefficients as well as evaluations.
On the unit disk it gives $p([c^\ell]Q(c))\le nP_*+D\log(nd+1)$.
Evaluating at $t^{-r_\infty}u$ can increase any coefficient price by
at most $nd r_\infty$. The price of their sum incurs another
$D\log(nd+1)$, proving $P_\infty-P_*\le d r_\infty$ in the limit.

To prove the other inequality, we use interpolation on each inner
disk to bound the low Taylor coefficients at its center and on the
outer disk to bound the high coefficients at infinity. More
precisely, interpolating $Q(c_j+t^{r_j}v)$ and $Q(t^{-r_\infty}v)$
in $v$ and removing the scalar power of $t$ gives
\begin{align}
 p([v^\ell]Q(c_j+v)) &\le nP_j+r_j\ell+D\log(nd+1),\notag\\
 p([c^\ell]Q(c)) &\le nP_\infty-r_\infty\ell+D\log(nd+1).
 \label{eq:valued-jets}
\end{align}

Suppose the sum of the inner drops per unit length exceeds the outer
growth per unit length. Choose $a_j$ strictly between zero and
$(P_*-P_j)/r_j$ on the arms with positive drop and zero on the other
arms, so that
\[
 \sum_j a_j>b>(P_\infty-P_*)/r_\infty
\]
for some $b$. Put $m=\lceil nb\rceil$ and $m_j=\lfloor na_j\rfloor$.
The strict inequalities and Equation~\eqref{eq:valued-jets} give
$\epsilon>0$ such that, for all large $n$, the coefficients
\[
 [c^\ell]Q(c)\quad(m\le\ell\le nd),\qquad
 [v^\ell]Q(c_j+v)\quad(0\le\ell<m_j)
\]
all have price at most $n(P_*-\epsilon)$. Together they determine $Q(u)$ by an integral linear combination.
Indeed, subtract the high-degree part of $Q$. The remaining polynomial
has degree $<m$ and its Taylor coefficients through order $m_j-1$
give its residue modulo $(c-c_j)^{m_j}$ at each center. For large
$n$, the total multiplicity $\sum_j m_j$ is at least $m$. The moduli
are pairwise comaximal over $\mathcal O$ because the center
differences are units. The Chinese remainder theorem thus recovers
the low-degree polynomial by an $\mathcal O$-linear map in monomial
coordinates. Computing the Taylor coefficients of the subtracted
part and evaluating at $u$ are integral operations as well. These
identities hold coefficientwise for tensors over $K$.

The reconstruction expresses $Q(u)$ using only $O(n+1)$ coefficients
from the displayed list. The sum bound would give
$nP_*=p(Q(u))\le n(P_*-\epsilon)+O(\log(n+1))$, a contradiction.

For the assisted cost, fix a power $m$ in
Equation~\eqref{eq:prices-beta} and write $A(c,\mathbf z)$ for the
closure of $T(c)^{\otimes m}$ with functional entries
$\mathbf z=(z_1,\ldots,z_s)$. These entries occur polynomially.
For any fixed $c$, the supremum of $p(A(c,\mathbf z))$ over
$\mathcal O^s$ equals the supremum over $F^s$. Indeed, interpolate
$A(c,\mathbf z)^{\otimes n}$ on a rectangular grid in $F^s$. For
fixed $m$, there are $O_m(n+1)$ nodes per coordinate and polynomially
many nodes in total. Evaluation at integral arguments has integral
weights, so the logarithmic cost per tensor copy tends to zero.

Consequently $\beta_p(T(c))$ is the supremum of the functions
$p(A(c,\mathbf z))/m$ over $m$ and constant closing data. Omit
identically zero tests when $T\ne0$. The other tests are nonzero at
all the generic points. Their specialization inequalities pass to
the supremum, and taking suprema over $a\in F$ and over tests in
either order gives the generic-value identity. Their weighted-average
inequalities pass to the supremum because each weight is positive.
Finally, each test satisfies the outer growth bound $d r_\infty$
after division by $m$, since $\deg_c A\le md$. This common bound
also holds for the supremum.
\end{proof}
